% !TeX root = ../main.tex
\section{Queue-Depth Sweep}\label{sec:queue-depth}

\begin{hypothesis}
  IOPS should rise steeply with queue depth as the SSD's parallel NAND channels fill, then saturate at a
  knee beyond which latency grows roughly linearly with queue depth for no extra throughput. By Little's
  Law the number of outstanding I/Os $L=\lambda W$ should track the configured iodepth while the queue
  stays full; where it stops tracking reveals a host-side or engine limit.
\end{hypothesis}

\begin{knobs}
  Fixed       & \qty{4}{\kibi\byte} random; one job; \code{O_DIRECT}; 5 repetitions of \qty{4}{\second}
                (writes also stop at \qty{256}{\mebi\byte} and start after a \qty{15}{\second} pause) \\
  Varied      & iodepth 1, 2, 4, \dots, 256; 100\% read and 100\% write \\
  Knee rule   & the smallest QD whose doubling raises median IOPS by less than \qty{10}{\percent}
                (computed by \code{run.sh} into \code{queue_depth_knee.csv}) \\
  Effective QD & checked in-band with Little's Law (\cref{tab:littles-law}); no \code{iostat} data \\
  Script      & \code{SSD_Performance_Profiling/scripts/run.sh queue_depth} \\
\end{knobs}

\pgfplotsset{qdplot/.style={acs, height=0.85\linewidth, xlabel={IOPS},
  ylabel={Latency (\si{\micro\second})}, legend pos=north west, legend style={font=\scriptsize},
  scaled x ticks=base 10:-3, xtick scale label code/.code={$\times10^{3}$},
  nodes near coords, point meta=explicit symbolic,
  every node near coord/.append style={font=\tiny, anchor=south east, inner sep=1pt}},
  qdbars/.style={error bars/.cd, x dir=both, x explicit, y dir=both, y explicit,
  error bar style={black!45}}}
\newcommand{\qdcurve}[2]{% {mode}{latency metric}
  \addplot+[qdbars] table[x=iops_median, y=#2_median, meta=qd, only if={mode}{#1},
    x error plus expr=\thisrow{iops_q75}-\thisrow{iops_median},
    x error minus expr=\thisrow{iops_median}-\thisrow{iops_q25},
    y error plus expr=\thisrow{#2_q75}-\thisrow{#2_median},
    y error minus expr=\thisrow{#2_median}-\thisrow{#2_q25}] {data/queue_depth.csv};}
\begin{figure}[!htbp]
  \centering
  \IfDataExists{queue_depth}{%
    \begin{subfigure}[t]{0.49\linewidth}
      \begin{tikzpicture}
        \begin{axis}[ymode=log, qdplot]
          \qdcurve{read}{lat_mean_us}\addlegendentry{mean}
          \qdcurve{read}{lat_p99_us}\addlegendentry{p99}
        \end{axis}
      \end{tikzpicture}
      \caption{4 KiB random read.}
    \end{subfigure}\hfill
    \begin{subfigure}[t]{0.49\linewidth}
      \csvget\ssdHasW{queue_depth}{mode=write}{iops_median}%
      \ifx\ssdHasW\empty\acsmissing{data/queue_depth.csv (write rows)}\else
      \begin{tikzpicture}
        \begin{axis}[ymode=log, qdplot]
          \qdcurve{write}{lat_mean_us}\addlegendentry{mean}
          \qdcurve{write}{lat_p99_us}\addlegendentry{p99}
        \end{axis}
      \end{tikzpicture}\fi
      \caption{4 KiB random write.}
    \end{subfigure}%
  }
  \caption{Throughput vs.\ latency as queue depth grows; each point is labelled with its QD. Bars show
  the IQR across repetitions in both IOPS and latency. Read IOPS reach
  \csvq{queue_depth_knee}{mode=read}{knee_iops} at the knee (QD~\csvq{queue_depth_knee}{mode=read}{knee_qd})
  and at most \csvq{queue_depth_knee}{mode=read}{max_iops} at QD=256; deeper queues then add only latency.
  In (b), the QD=4 point at \csvq{queue_depth}{qd=4,mode=write}{iops_median} IOPS is a cache-exhaustion
  stall (see text).}\label{fig:queue-depth}
\end{figure}

\subsection{Knee}
We define the knee as \emph{the smallest QD whose doubling raises median IOPS by less than
\qty{10}{\percent}}: past it, most extra concurrency becomes latency rather than throughput.

\textbf{Reads.} The rule gives QD~\csvq{queue_depth_knee}{mode=read}{knee_qd}
(\csvq{queue_depth_knee}{mode=read}{knee_iops} IOPS,
\csvq[0]{queue_depth_knee}{mode=read}{knee_lat_us}\,\si{\micro\second} mean). Doubling to
QD~\csvq{queue_depth_knee}{mode=read}{next_qd} adds only
\csvq[1]{queue_depth_knee}{mode=read}{gain_pct}\,\% IOPS, while QD 32$\to$64 still gained
\qty{12}{\percent}, so the knee lies between 32 and 64.

\textbf{Writes.} The rule returns QD~\csvq{queue_depth_knee}{mode=write}{knee_qd}
(\csvq{queue_depth_knee}{mode=write}{knee_iops} IOPS), because the next point collapsed: at QD=4 IOPS
change by \csvq[1]{queue_depth_knee}{mode=write}{gain_pct}\,\%, to
\csvq{queue_depth}{qd=4,mode=write}{iops_median}. All five QD=4 repetitions ran at
\qtyrange{61}{67}{\mebi\byte\per\second}, the post-cliff QLC rate of \cref{sec:device-state}. Their median
latency stayed at \csvq[1]{queue_depth}{qd=4,mode=write}{lat_p50_us_median}\,\si{\micro\second} and p99
at \csvq[0]{queue_depth}{qd=4,mode=write}{lat_p99_us_median}\,\si{\micro\second}, but under
\qty{1}{\percent} of writes waited \ensuremath{\sim}\qty{33}{\milli\second} (p99.9) for the buffer to drain, which sets
the mean of \csvq[0]{queue_depth}{qd=4,mode=write}{lat_mean_us_median}\,\si{\micro\second}. $L/\text{QD}$
stays near 1 (\cref{tab:littles-law}), so the queue was full and the drive simply accepted writes no
faster. The point came after \qty{2.5}{\gibi\byte} of writes in this step and \qty{24}{\gibi\byte} in
\code{block_size}, with the cache not fully recovered from the \qty{60}{\gibi\byte} fill; the first QD=8
repetition was still slow (\num{20076} IOPS), the other four ran at \num{223}k. Ignoring the stalled
point, IOPS rise from \csvq{queue_depth}{qd=1,mode=write}{iops_median} at QD=1 to
\csvq{queue_depth}{qd=16,mode=write}{iops_median} at QD=16, the 10\,\% rule would place the knee at
QD=16 (QD 8$\to$16 gains \qty{13}{\percent}, 16$\to$32 \qty{-2.5}{\percent}), and the curve stays at
\csvq{queue_depth}{qd=32,mode=write}{iops_median}--\csvq{queue_depth}{qd=128,mode=write}{iops_median} up
to QD=256. The rule's QD=2 is thus an artifact of the stall; in-cache writes saturate at QD\,$\approx$\,16.
The plateau of \ensuremath{\sim}\num{250}k IOPS $\times$ \qty{4}{\kibi\byte} $\approx$
\qty{0.95}{\gibi\byte\per\second} is the SLC ingest rate: writes saturate the buffer's intake, not the
NAND parallelism that limits reads.

\subsection{Little's Law as a consistency check}
Little's Law \cite{little1961proof} states $L=\lambda W$ for the mean number in system $L$, throughput
$\lambda$ and mean time in system $W$. In a closed loop it holds by construction:
$\lambda W = \text{QD}\cdot W/(W+g)$, where $g$ is the host gap between reaping and resubmitting. So
$L/\text{QD}\approx1$ shows that the timers are consistent and the host keeps up ($g\ll W$), while a
clearly lower ratio would expose a host bottleneck. As a control, \code{make test} on tmpfs, where I/O
completes inside \code{io_submit}, measured $L/\text{QD}\approx0.57$ ($L=4.56$ at QD=8;
\code{data/raw/validate_tmpfs/20261021-231433.txt}), which the harness flags. $L$ includes time in the
kernel block layer, so device-queue occupancy can only be lower.

\begin{table}[!htbp]
  \centering\small
  \caption{Little's Law per queue depth: $\lambda$ (IOPS) and $W$ (mean latency, \si{\micro\second}) are
  medians with the IQR in brackets; $L=\lambda W$ is computed per repetition and then summarized.
  \code{4 KiB} random.}\label{tab:littles-law}
  \IfDataExists{queue_depth}{%
    \ssdload{queue_depth}\def\ssdRows{}%
    \pgfplotsforeachungrouped \ssdR in {0,...,\ssdLast}{%
      \ssdstr{mode}\ssdM\ssdstr{qd}\ssdQD
      \ssdiqr{iops}\ssdA\ssdiqr[1]{lat_mean_us}\ssdB\ssdiqr[2]{littles_l}\ssdC
      \pgfplotstablegetelem{\ssdR}{littles_l_median}\of\ssdT\let\ssdLraw\pgfplotsretval
      \pgfkeys{/pgf/fpu=true}\pgfmathparse{\ssdLraw/\ssdQD}%
      \pgfmathprintnumberto[fixed,precision=3]{\pgfmathresult}{\ssdD}\pgfkeys{/pgf/fpu=false}%
      \eappto\ssdRows{\expandonce\ssdM & \ssdQD & \expandonce\ssdA & \expandonce\ssdB & \expandonce\ssdC &
        \expandonce\ssdD \noexpand\\}}%
    \begin{tabular}{@{}lrrrrr@{}}
      \toprule
      Mode & QD & $\lambda$ (IOPS) [IQR] & $W$ [IQR] & $L=\lambda W$ [IQR] & $L/\text{QD}$ \\
      \midrule
      \ssdRows
      \bottomrule
    \end{tabular}}
\end{table}

\paragraph{Inference: internal parallelism.} Applied to the device, with
$\lambda_{\max}\approx\csvq{queue_depth_knee}{mode=read}{max_iops}$ reads/s and an assumed service time $S$
between \qty{40}{\micro\second} and the QD=1 latency
(\csvq[0]{queue_depth}{qd=1,mode=read}{lat_mean_us_median}\,\si{\micro\second}), only
$L_\text{busy}=\lambda_{\max} S\approx 6$--$8$ reads are in service at once; of the 64--256 outstanding
past the knee, the rest wait in queues. With 1-Tbit dies (4 in a \qty{512}{\giga\byte} drive; Micron
does not publish this) that is about two reads per die, a plausibility argument only.

\subsection{Analysis}
Read IOPS climb from \csvq{queue_depth}{qd=1,mode=read}{iops_median} at QD=1 to
\csvq{queue_depth}{qd=16,mode=read}{iops_median} at QD=16 and
\csvq{queue_depth}{qd=32,mode=read}{iops_median} at QD=32, then flatten at
\csvq{queue_depth}{qd=64,mode=read}{iops_median}--\csvq{queue_depth}{qd=256,mode=read}{iops_median}.
Scaling is sub-linear from the start. Beyond the knee mean latency doubles with each doubling of QD
(\csvq[0]{queue_depth}{qd=64,mode=read}{lat_mean_us_median}$\to$%
\csvq[0]{queue_depth}{qd=128,mode=read}{lat_mean_us_median}$\to$%
\csvq[0]{queue_depth}{qd=256,mode=read}{lat_mean_us_median}\,\si{\micro\second}), the textbook
saturation signature: with $\lambda$ flat, every added request simply waits.

For reads $L/\text{QD}$ stays at 0.987--0.999 up to QD=256 (\cref{tab:littles-law}), so the single
\code{libaio} thread keeps the queue full and the saturation is the device's. For writes it is
0.990--0.998, including the stalled QD=4 point.

In practice, reads need QD\,$\approx$\,32--64 for nearly all the throughput at a quarter to half a
millisecond; deeper queues only add latency. Writes need only QD\,$\approx$\,16 while the SLC cache has
room; once it is full, QLC programming sets the rate at any depth (\cref{sec:device-state}).
